% GENERATED FILE. Edit canonical lessons, then run npm run generate:books.
% canonical-source-hash: fnv1a64:c24ac187873d8e90
% canonical-lessons: BN-C458-sot, BN-C458-kopot, BN-C458-bishvosto, BN-C458-dhurto, BN-C458-thokano

\chapter{Honest and Deceitful}
\label{ch:bn-honest-and-deceitful}
\hlchaptermodality{458}

\begin{chapteropening}
\textbf{By the end of this chapter:} \emph{I can say whether someone is honest or a cheat.}

Complete the last lesson of chapter 458: I can say whether someone is honest or a cheat.
\end{chapteropening}

\section[\texorpdfstring{\bn{সৎ} (sôt)}{সৎ (sôt)}]{\bn{সৎ} (sôt) — honest; upright}

\label{lesson:BN-C458-sot}



\begin{quote}
Before the new one: say the Bengali for helpful to others, then the Bengali for noble.
\end{quote}

\subsection*{You'll want to know: \bn{সৎ}}
\textbf{\bn{সৎ}} — \emph{sôt} — \textquotedblleft{}honest; upright\textquotedblright{}.

\begin{quote}
\bn{সৎ} \bn{লোকটা} \bn{মানিব্যাগটা} \bn{পুলিশকে} \bn{দিল।}
\end{quote}

\begin{quote}
\emph{sôt lokṭā mānibyāgṭā pulishke dilô.}
\end{quote}

\begin{quote}
The honest man handed the wallet to the police.
\end{quote}

\bn{সৎ} also marks a step-relation: \bn{সৎ} \bn{ভাই}, a half-brother.

\begin{cousinweb}
\textbf{\bn{সৎ}} (\emph{sôt}) is a tatsama, Sanskrit \emph{sat}, \textquotedblleft{}being, existing\textquotedblright{} and \textquotedblleft{}true, good, honest\textquotedblright{} (Monier-Williams); it ends in the new letter \bn{ৎ}.
\end{cousinweb}

\begin{grammarlens}[title={What you've built}]
One more word for describing what you lived and felt.
\end{grammarlens}

\subsection*{Your turn}
\begin{itemize}
\raggedright
  \item \emph{Say it:} \emph{sôt}
  \item \emph{Say it:} \emph{sôt}, once more
  \item \emph{Say it:} \emph{sôt lokṭā mānibyāgṭā pulishke dilô.} — The honest man handed the wallet to the police.
  \item \emph{Recall:} say the Bengali for helpful to others, then the Bengali for noble, then say \emph{sôt} again
\end{itemize}

\begin{tcolorbox}[breakable,colback=teal!4,colframe=teal!35!black,title={Before you move on}]
What does \bn{সৎ} mean? (\textquotedblleft{}Honest; upright\textquotedblright{}.) Say it once more.
\end{tcolorbox}
\section[\texorpdfstring{\bn{কপট} (kôpôṭ)}{কপট (kôpôṭ)}]{\bn{কপট} (kôpôṭ) — insincere; two-faced; feigned}

\label{lesson:BN-C458-kopot}



\begin{quote}
Before the new one: say the Bengali for noble, then the Bengali for honest.
\end{quote}

\subsection*{You'll want to know: \bn{কপট}}
\textbf{\bn{কপট}} — \emph{kôpôṭ} — \textquotedblleft{}insincere; two-faced; feigned\textquotedblright{}.

\begin{quote}
\bn{লোকটার} \bn{কপট} \bn{হাসি} \bn{দেখে} \bn{মীরা} \bn{তাকে} \bn{বিশ্বাস} \bn{করল} \bn{না।}
\end{quote}

\begin{quote}
\emph{lokṭār kôpôṭ hāsi dekhe mīrā tāke bishvās kôrlô nā.}
\end{quote}

\begin{quote}
Seeing the man's false smile, Mira did not trust him.
\end{quote}

\bn{কপট} \bn{রাগ}, pretended anger: the kind a grandmother puts on.

\begin{cousinweb}
\textbf{\bn{কপট}} (\emph{kôpôṭ}) is a tatsama, Sanskrit \emph{kapaṭa}, a noun for fraud and deceit (Monier-Williams); Bengali uses it as an adjective, for a smile or a feeling put on to mislead. It is a different word from Sanskrit \emph{kapāṭa}, a door leaf.
\end{cousinweb}

\begin{grammarlens}[title={What you've built}]
One more word for describing what you lived and felt.
\end{grammarlens}

\subsection*{Your turn}
\begin{itemize}
\raggedright
  \item \emph{Say it:} \emph{kôpôṭ}
  \item \emph{Say it:} \emph{kôpôṭ}, once more
  \item \emph{Say it:} \emph{lokṭār kôpôṭ hāsi dekhe mīrā tāke bishvās kôrlô nā.} — Seeing the man's false smile, Mira did not trust him.
  \item \emph{Recall:} say the Bengali for noble, then the Bengali for honest, then say \emph{kôpôṭ} again
\end{itemize}

\begin{tcolorbox}[breakable,colback=teal!4,colframe=teal!35!black,title={Before you move on}]
What does \bn{কপট} mean? (\textquotedblleft{}Insincere; two-faced; feigned\textquotedblright{}.) Say it once more.
\end{tcolorbox}
\section[\texorpdfstring{\bn{বিশ্বস্ত} (bishvôstô)}{বিশ্বস্ত (bishvôstô)}]{\bn{বিশ্বস্ত} (bishvôstô) — trustworthy; loyal}

\label{lesson:BN-C458-bishvosto}



\begin{quote}
Before the new one: say the Bengali for honest, then the Bengali for two-faced.
\end{quote}

\subsection*{You'll want to know: \bn{বিশ্বস্ত}}
\textbf{\bn{বিশ্বস্ত}} — \emph{bishvôstô} — \textquotedblleft{}trustworthy; loyal\textquotedblright{}.

\begin{quote}
\bn{আমাদের} \bn{পুরনো} \bn{দারোয়ান} \bn{বিশ্বস্ত} \bn{লোক।}
\end{quote}

\begin{quote}
\emph{āmāder purôno dāroyān bishvôstô lok.}
\end{quote}

\begin{quote}
Our old gatekeeper is a trustworthy man.
\end{quote}

\begin{cousinweb}
\textbf{\bn{বিশ্বস্ত}} (\emph{bishvôstô}) is a tatsama, Sanskrit \emph{viśvasta}, \textquotedblleft{}trusted, faithful\textquotedblright{} (Monier-Williams), from the same root as \bn{বিশ্বাস} \bn{করা} (taught), to trust.
\end{cousinweb}

\begin{grammarlens}[title={What you've built}]
One more word for describing what you lived and felt.
\end{grammarlens}

\subsection*{Your turn}
\begin{itemize}
\raggedright
  \item \emph{Say it:} \emph{bishvôstô}
  \item \emph{Say it:} \emph{bishvôstô}, once more
  \item \emph{Say it:} \emph{āmāder purôno dāroyān bishvôstô lok.} — Our old gatekeeper is a trustworthy man.
  \item \emph{Recall:} say the Bengali for honest, then the Bengali for two-faced, then say \emph{bishvôstô} again
\end{itemize}

\begin{tcolorbox}[breakable,colback=teal!4,colframe=teal!35!black,title={Before you move on}]
What does \bn{বিশ্বস্ত} mean? (\textquotedblleft{}Trustworthy; loyal\textquotedblright{}.) Say it once more.
\end{tcolorbox}
\section[\texorpdfstring{\bn{ধূর্ত} (dhūrtô)}{ধূর্ত (dhūrtô)}]{\bn{ধূর্ত} (dhūrtô) — cunning; crafty}

\label{lesson:BN-C458-dhurto}



\begin{quote}
Before the new one: say the Bengali for two-faced, then the Bengali for trustworthy.
\end{quote}

\subsection*{You'll want to know: \bn{ধূর্ত}}
\textbf{\bn{ধূর্ত}} — \emph{dhūrtô} — \textquotedblleft{}cunning; crafty\textquotedblright{}.

\begin{quote}
\bn{ধূর্ত} \bn{শিয়ালটা} \bn{কাকের} \bn{রুটি} \bn{নিয়ে} \bn{পালাল।}
\end{quote}

\begin{quote}
\emph{dhūrtô shiyālṭā kāker ruṭi niye pālālô.}
\end{quote}

\begin{quote}
The crafty jackal ran off with the crow's bread.
\end{quote}

\begin{cousinweb}
\textbf{\bn{ধূর্ত}} (\emph{dhūrtô}) is a tatsama, Sanskrit \emph{dhūrta}, \textquotedblleft{}cunning, crafty\textquotedblright{}, and as a noun \textquotedblleft{}a rogue, a cheat\textquotedblright{} (Monier-Williams).
\end{cousinweb}

\begin{grammarlens}[title={What you've built}]
One more word for describing what you lived and felt.
\end{grammarlens}

\subsection*{Your turn}
\begin{itemize}
\raggedright
  \item \emph{Say it:} \emph{dhūrtô}
  \item \emph{Say it:} \emph{dhūrtô}, once more
  \item \emph{Say it:} \emph{dhūrtô shiyālṭā kāker ruṭi niye pālālô.} — The crafty jackal ran off with the crow's bread.
  \item \emph{Recall:} say the Bengali for two-faced, then the Bengali for trustworthy, then say \emph{dhūrtô} again
\end{itemize}

\begin{tcolorbox}[breakable,colback=teal!4,colframe=teal!35!black,title={Before you move on}]
What does \bn{ধূর্ত} mean? (\textquotedblleft{}Cunning; crafty\textquotedblright{}.) Say it once more.
\end{tcolorbox}
\section[\texorpdfstring{\bn{ঠকানো} (ṭhôkāno)}{ঠকানো (ṭhôkāno)}]{\bn{ঠকানো} (ṭhôkāno) — to cheat; to swindle}

\label{lesson:BN-C458-thokano}



\begin{quote}
Before the new one: say the Bengali for trustworthy, then the Bengali for cunning.
\end{quote}

\subsection*{You'll want to know: \bn{ঠকানো}}
\textbf{\bn{ঠকানো}} — \emph{ṭhôkāno} — \textquotedblleft{}to cheat; to swindle\textquotedblright{}.

\begin{quote}
\bn{দোকানদার} \bn{আমাকে} \bn{দশ} \bn{টাকা} \bn{ঠকাল।}
\end{quote}

\begin{quote}
\emph{dokāndār āmāke dôsh ṭākā ṭhôkālô.}
\end{quote}

\begin{quote}
The shopkeeper cheated me out of ten taka.
\end{quote}

Past (he/she …): \bn{তিনি} \bn{খদ্দেরকে} \bn{ঠকালেন}, \emph{tini khôdderke ṭhôkālen}.

\begin{cousinweb}
\textbf{\bn{ঠকানো}} (\emph{ṭhôkāno}) is listed by CDIAL for Bengali, as deceiving, under a reconstructed \emph{ṭhagg}, \textquotedblleft{}cheat\textquotedblright{} (CDIAL 5489); \emph{ṭhôkā}, to be cheated, is its plain partner.
\end{cousinweb}

\begin{grammarlens}[title={What you've built}]
One more word for describing what you lived and felt.
\end{grammarlens}

\subsection*{Your turn}
\begin{itemize}
\raggedright
  \item \emph{Say it:} \emph{ṭhôkāno}
  \item \emph{Say it:} \emph{ṭhôkāno}, once more
  \item \emph{Say it:} \emph{dokāndār āmāke dôsh ṭākā ṭhôkālô.} — The shopkeeper cheated me out of ten taka.
  \item \emph{Recall:} say the Bengali for trustworthy, then the Bengali for cunning, then say \emph{ṭhôkāno} again
\end{itemize}

\begin{tcolorbox}[breakable,colback=teal!4,colframe=teal!35!black,title={Before you move on}]
What does \bn{ঠকানো} mean? (\textquotedblleft{}To cheat; to swindle\textquotedblright{}.) Say it once more.
\end{tcolorbox}
